\documentclass[12pt]{article}
\usepackage[utf8]{inputenc}
\usepackage[T1]{fontenc}
\usepackage[margin=1in]{geometry}
\usepackage{setspace}
\usepackage{xcolor}
\usepackage[colorlinks=true, urlcolor=blue!60!black, linkcolor=blue!60!black]{hyperref}
\setlength{\parindent}{2em}
\setlength{\parskip}{1.25em}
\renewcommand{\baselinestretch}{1.0}
\begin{document}
\begin{center}
{\LARGE \textit{Cooper v. Aaron}, 358 U.S. 1 (1958)\par}
\vspace{0.8em}
{\large Case Brief --- Four-Page Edition\par}
\end{center}
\vspace{1em}

\section{Memory Jogger}

Facing Arkansas's open defiance of \textit{Brown}, all nine Justices signed an opinion holding that disorder provoked by state officials cannot justify suspending a desegregation plan, and declaring that the Court's interpretation of the Fourteenth Amendment is ``the supreme law of the land'' binding every state official (\href{https://supreme.justia.com/cases/federal/us/358/1/}{Justia}).

\section{Detailed Case Facts}

Days after \textit{Brown v. Board of Education} (1954), the Little Rock School Board announced it would comply, and in May 1955 it adopted a gradual plan to begin at Central High School in fall 1957 and finish by 1963. Federal courts upheld the plan. Meanwhile Arkansas launched ``Massive Resistance'': a 1956 state constitutional amendment commanded the legislature to oppose ``in every Constitutional manner the Unconstitutional desegregation decisions'' of the Court, and 1957 statutes excused children from attending integrated schools.

On September 2, 1957, the night before the ``Little Rock Nine'' were to enter Central High, Governor Orval Faubus, without any request from local authorities or prior violence, deployed the Arkansas National Guard to place the school ``off limits'' to Black students. After a federal court enjoined him, the Guard withdrew, but a hostile mob forced police to remove the students. President Eisenhower sent U.S. Army troops, and eight of the nine students finished the year.

\section{Procedural History}

In February 1958, citing ``chaos, bedlam and turmoil,'' the School Board asked the district court to suspend the plan for two and one-half years. The court granted the suspension; the Eighth Circuit reversed but stayed its mandate. In a special August Term, the Supreme Court granted certiorari on September 11, affirmed per curiam on September 12, three days before school opened, and filed its full opinion on September 29, 1958.

\section{Judicial Votes}

Affirmed, \textbf{9--0}.
\begin{itemize}
\item \textbf{Opinion of the Court:} jointly signed by Chief Justice Warren and Justices Black, Frankfurter, Douglas, Burton, Clark, Harlan, Brennan, and Whittaker, the only time every Justice personally signed the Court's opinion.
\item \textbf{Concurrence:} Justice Frankfurter (October 6, 1958).
\item \textbf{Dissents:} None.
\end{itemize}

\section{Holding}

A good-faith school board may not suspend a court-approved desegregation plan because of disorder provoked by the Governor and Legislature, and state officials are bound by the Court's interpretation of the Fourteenth Amendment. School officials act as agents of the State and cannot invoke their own good faith, or disorder created by other officials, to suspend constitutional rights. Under \textit{Marbury v. Madison} and Article VI, the federal judiciary is supreme in the exposition of the Constitution, so \textit{Brown} is the ``supreme Law of the Land'' binding all state officials.

\section{Analysis of Opinions}

\textbf{Joint opinion: state action.} The Court accepted the Board's good faith and the findings of turmoil, but the conditions were ``directly traceable'' to the Governor and Legislature. Under \textit{Ex parte Virginia}, a State acts through all its officers, so the Board could not plead obstruction by other state officials. Under \textit{Buchanan v. Warley}, ``law and order are not here to be preserved by depriving the Negro children of their constitutional rights.''

\textbf{Joint opinion: judicial supremacy.} Though state action sufficed, the Court answered Arkansas's claim that non-parties were not bound by \textit{Brown}. Reading the Supremacy Clause with \textit{Marbury}, the Court held that its \emph{interpretation} in \textit{Brown} binds legislators, executives, and judges through their Article VI oath. Citing \textit{Ableman v. Booth}, \textit{United States v. Peters}, and \textit{Sterling v. Constantin}, it noted that the three Justices appointed since \textit{Brown} joined in reaffirming it.

\textbf{Frankfurter (concurring).} Peaceful accommodation had been underway before the State used the militia to thwart law. Yielding would ``enthrone official lawlessness'' and mean ``that law should bow to force.'' Orderly criticism and legal change differ from defiance.

\textbf{Later developments.} \textit{United States v. Nixon} (1974) and \textit{City of Boerne v. Flores} (1997) extended \textit{Cooper}'s supremacy framework to the President and Congress.

\section{Examples: Future Applications}

\begin{enumerate}
\item \textbf{Same-Side: A Non-Party Governor.} A Governor orders prosecutors to keep enforcing a statute identical to one the Court struck down in another State's case. Under \textit{Cooper}, that violates Article VI; a departmentalist would distinguish defying a judgment from declining to treat a precedent as settling his own judgment.
\item \textbf{Same-Side: A County Clerk.} A clerk stops issuing marriage licenses after the Court requires same-sex marriage licensing. She acts for the State, is bound through her oath, and may be held in contempt.
\item \textbf{Opposite-Side: A Presidential Veto.} A President vetoes a bank charter the Court would uphold. The veto defies no judgment and denies no declared right, as in Jackson's 1832 veto.
\item \textbf{Opposite-Side: Another State's District Court.} State officials keep enforcing a voter-ID law like one a district court elsewhere enjoined. \textit{Cooper} attaches to the Supreme Court's exposition, not every lower-court opinion.
\item \textbf{Fence-Sitter: A Test-Case Statute.} A legislature enacts a law conflicting with precedent to invite overruling, like the Mississippi law in \textit{Dobbs}. It resembles an ``evasive scheme,'' yet Frankfurter recognized orderly efforts to change the law.
\end{enumerate}

\section{Critique}

\textbf{Textual slippage.} Article VI lists the Constitution, federal laws, and treaties as ``the supreme Law of the Land,'' not Supreme Court opinions. Critics such as Edwin Meese III and departmentalist scholars argue that \textit{Cooper} conflates the Constitution with the Court's fallible interpretations and leaps from \textit{Marbury}'s case-specific review to universal interpretive supremacy.

\textbf{The Lincoln objection.} Taken literally, \textit{Cooper} would mean that Lincoln violated his oath by refusing to treat \textit{Dred Scott} as a permanent political rule. Even Frankfurter conceded that ``[e]ven this Court has the last say only for a time.''

\section{Key Quotations}

\begin{quote}
``This decision declared the basic principle that the federal judiciary is supreme in the exposition of the law of the Constitution, and that principle has ever since been respected by this Court and the Country as a permanent and indispensable feature of our constitutional system. It follows that the interpretation of the Fourteenth Amendment enunciated by this Court in the Brown case is the supreme law of the land, and Art.\ VI of the Constitution makes it of binding effect on the States `any Thing in the Constitution or Laws of any State to the Contrary notwithstanding.'\thinspace''
\end{quote}
\noindent --- Opinion of the Court, 358 U.S. at 18.

\begin{quote}
``In short, the constitutional rights of children not to be discriminated against in school admission on grounds of race or color declared by this Court in the Brown case can neither be nullified openly and directly by state legislators or state executive or judicial officers nor nullified indirectly by them through evasive schemes for segregation whether attempted `ingeniously or ingenuously.'\thinspace''
\end{quote}
\noindent --- Opinion of the Court, 358 U.S. at 17.

\begin{quote}
``The constitutional rights of respondents are not to be sacrificed or yielded to the violence and disorder which have followed upon the actions of the Governor and Legislature.''
\end{quote}
\noindent --- Opinion of the Court, 358 U.S. at 16.

\end{document}
